\documentclass[11pt]{article}
\usepackage[a4paper,margin=0.85in]{geometry}
\usepackage{amsmath,amssymb}
\usepackage[hidelinks]{hyperref}
\title{Literature status: general relaxation in $BD$ and recession functions}
\author{Source arXiv:1008.2089; answers arXiv:1903.05771 and 1911.01356}
\date{13 August 2026}
\begin{document}
\maketitle
\noindent\textbf{Status: literature already answered.}

\section*{Original question}
Remark 6.9 of Rindler \cite{Rindler} (source PDF p.~44) asks for a
relaxation theorem for integrands $f$ that are not symmetric-quasiconvex,
with $f$ replaced by its symmetric-quasiconvex envelope $SQf$.  It records
two obstacles: whether the strong recession function $(SQf)^\infty$ exists
was unknown, and the paper's framework could not establish lower
semicontinuity after replacing it by the generalized recession function
$(SQf)^\#$.

\section*{Exact later resolution}
Kosiba--Rindler \cite{KR}, Theorem 1.3 (supporting PDF pp.~3--4), proves the
desired weak-star relaxation in the standard homogeneous setting.  If
$f\colon\mathbb R^{d\times d}_{\rm sym}\to[0,\infty)$ is continuous and
\[
  m|A|\leq f(A)\leq M(1+|A|),
\]
then the relaxation from $LD(\Omega)$ to $BD(\Omega)$ is
\[
 \overline F[u,\Omega]
 =\int_\Omega (SQf)(Eu)\,dx
 +\int_\Omega (SQf)^\#\!\left(\frac{dE^s u}{d|E^s u|}\right)d|E^s u|.
\]
Immediately afterward, the authors emphasize that this theorem requires no
assumption on a recession function and explicitly cite Rindler's source paper
as the preceding strong-recession result and discussion.

The later survey of De Philippis--Rindler \cite{DPR} makes the logical status
fully explicit.  On supporting PDF pp.~25--26 it says that, in general, the
strong recession function $(SQf)^\infty$ does not exist (citing a published
counterexample), introduces $(SQf)^\#$ for precisely this reason, and states
the same result as Theorem 4.5.  Thus the source's existence question has a
negative answer, while its relaxation objective is achieved by the proposed
generalized-recession formula.

\section*{Provenance, scope, and search evidence}
This is not an agent-derived implication: \cite{KR} cites the source as its
reference [34], contrasts its Theorem 1.3 with that paper's result, and points
back to the source's recession-function remarks.  The survey cites the source
as [64] and the relaxation theorem as [51].  The theorem quoted above assumes
an $x$-independent continuous integrand with two-sided linear growth; it does
not by itself settle arbitrary dependence on $x$ or $u$.  It does exactly
resolve the homogeneous relaxation problem discussed in Remark 6.9, and the
failure of $(SQf)^\infty$ is general.

\begin{thebibliography}{9}
\bibitem{Rindler} F. Rindler,
\emph{Lower semicontinuity for integral functionals in the space of functions
of bounded deformation via rigidity and Young measures},
Arch. Ration. Mech. Anal. 202 (2011), 63--113; arXiv:1008.2089.
\bibitem{KR} K. Kosiba and F. Rindler,
\emph{On the relaxation of integral functionals depending on the symmetrized
gradient}, arXiv:1903.05771 (2020 version).
\bibitem{DPR} G. De Philippis and F. Rindler,
\emph{Fine properties of functions of bounded deformation---an approach via
linear PDEs}, arXiv:1911.01356 (2019).
\end{thebibliography}
\end{document}

